\section{总体相似度概览}

\begin{frame}{总体评分：LLM 与人类基线的相似度}
  \begin{block}{读图方式}
    SRS 是社会理性综合分，整合行为对齐 BAS 与价值对齐 VAS。
    分数越高，表示模型在当前两类社会决策范式中越接近人类基线。
  \end{block}

  \vspace{-0.4em}
  \begin{tikzpicture}
    \begin{axis}[
      xbar,
      width=0.97\textwidth,
      height=0.70\textheight,
      xmin=45,
      xmax=63,
      xlabel={SRS},
      ytick={1,...,19},
      yticklabels={
        Q3.6-max,Q3.6-flash,Q3.6-35B,DS-v4-flash T,GLM-5.1 NT,
        Q3.5-122B,Q3.6-plus NT,Q3.6-plus T,Q3.6-plus,DS-v3.2 T,
        Q3.5-35B,GLM-5.1 T,Q3.5-flash,Q3.5-plus,Q3.5-plus T,
        Q3.5-plus NT,Q3.5-397B,DS-v4-pro T,DS-v4-pro NT},
      y dir=reverse,
      tick label style={font=\scriptsize},
      yticklabel style={font=\tiny},
      label style={font=\small},
      xmajorgrids,
      grid style={BenchGray!18},
      nodes near coords,
      nodes near coords style={font=\tiny, anchor=west},
      point meta=x,
      every node near coord/.append style={
        /pgf/number format/fixed,
        /pgf/number format/precision=1
      },
    ]
      \addplot[fill=BenchBlue!75, draw=BenchBlue] coordinates {
        (48.6,1) (52.8,2) (53.0,3) (54.8,4) (56.5,5)
        (56.6,6) (57.1,7) (57.3,8) (57.4,9) (57.5,10)
        (57.8,11) (58.1,12) (58.3,13) (58.3,14) (58.6,15)
        (58.7,16) (58.8,17) (59.1,18) (61.4,19)
      };
    \end{axis}
  \end{tikzpicture}

  \vspace{-0.5em}
  \smallnote{T = 最大思考模式；NT = 非思考模式。}
\end{frame}

\begin{frame}{总体评分拆解：BAS 与 VAS 并不总是同步}
  \begin{columns}[T,totalwidth=\textwidth]
    \begin{column}{0.66\textwidth}
      \begin{tikzpicture}
        \begin{axis}[
          width=\textwidth,
          height=0.72\textheight,
          xmin=48,
          xmax=64,
          ymin=38,
          ymax=67,
          xlabel={BAS：行为对齐},
          ylabel={VAS：价值对齐},
          grid=both,
          grid style={BenchGray!15},
          legend style={at={(0.02,0.98)},anchor=north west,font=\scriptsize},
          tick label style={font=\scriptsize},
          label style={font=\small},
        ]
          \addplot[dashed, BenchGray, domain=48:64] {x};
          \addplot+[only marks, mark=*, mark size=2pt, color=BenchGreen] coordinates {
            (50.8,64.2) (59.4,50.2) (58.4,59.7) (57.5,65.4)
          };
          \addplot+[only marks, mark=square*, mark size=2pt, color=BenchOrange] coordinates {
            (55.0,61.3) (55.2,57.9)
          };
          \addplot+[only marks, mark=triangle*, mark size=2pt, color=BenchBlue] coordinates {
            (49.9,63.3) (51.8,63.8) (51.8,65.7) (52.3,64.3) (51.4,65.2)
            (51.6,65.5) (62.3,55.1) (58.6,47.5) (58.3,47.2) (57.3,39.9)
            (59.5,55.3) (59.9,54.8) (60.9,53.3)
          };
          \node[font=\scriptsize, anchor=west, text=BenchRed] at (axis cs:57.3,39.9) {Q3.6-max};
          \node[font=\scriptsize, anchor=south, text=BenchGreen] at (axis cs:57.5,65.4) {DS-v4-pro NT};
          \legend{BAS=VAS,DeepSeek,GLM,Qwen}
        \end{axis}
      \end{tikzpicture}
    \end{column}
    \begin{column}{0.31\textwidth}
      \begin{block}{观察}
        \begin{itemize}
          \item 总体 SRS 范围为 48.6--61.4。
          \item 一些模型 BAS 较高但 VAS 较低，说明具体选择接近不代表趋势结构接近。
          \item 因此后续结果分开看行为对齐与价值对齐。
        \end{itemize}
      \end{block}
    \end{column}
  \end{columns}
\end{frame}
